\subsection{Truth Tables (Cont.)}

\content{
        \begin{example} A friend tells you ``If you upload your picture to
        Facebook, you'll lose your job." Under what conditions is your friend
        wrong?
        \end{example}
}{% presentation
\vspace{2in}
}{% nopresentation
\vspace{1.5in}
}{% comments
Define events $p$ and $q$, go through all possibilities of true and false (e.g.
you upload the picture and you lose your job)
}

\content{
        \begin{example} Draw a truth table for the statement 
        $\neg (p \land q)$.
        \end{example}
}{% presentation
\vspace{1.5in}
}{% nopresentation
\vspace{1in}
}{% comments
}

\content{
        \begin{definition} (Related Statements) For any conditional, there are
        three related statements: 
        \begin{enumerate}
        \item (Converse) if $q$, then $p$ ($q\Rightarrow p$)
        \item (Inverse) if not $p$, then not $q$ ($\neg p \Rightarrow \neg q$)
        \item (Contrapositive) if not $q$, then not $p$ ($\neg q \Rightarrow
        \neg p$)
        \end{enumerate}
        \end{definition}
}{% presentation
}{% nopresentation
}{% comments
}

\content{
        \begin{example} Consider the statement ``If it is raining, then there
        are clouds in the sky". What are the converse, inverse, and
        contrapositive of this statement? Which are true?
        \end{example}
}{% presentation
\vspace{1.5in}
}{% nopresentation
\vspace{1.5in}
}{% comments
}

\content{
        \begin{example} Create a truth table for the conditional, converse,
        inverse, and contrapositive. Which are logically equivalent?
        \end{example}
}{% presentation
\vspace{2in}
}{% nopresentation
\vspace{1.5in}
}{% comments
}

\content{
        \begin{example} Consider the statement ``If you cook salmon in the break
        room oven, I will be mad at you". What are the converse, inverse, and
        contrapositive statements?
        \end{example}
}{% presentation
\vspace{2in}
}{% nopresentation
\vspace{1.5in}
}{% comments
}


\content{
        \begin{definition} (Biconditional) The biconditional statement is used
        when two statments are interchangable. Denoted $p\Leftrightarrow q$.

        The truth table of the biconditional is as follows: 

        \begin{tabular}{|c|c|c|}
        \hline
        $p$ & $q$ & $p\Leftrightarrow q$  \\
        \hline
        T & T & T \\
        \hline
        T & F & F \\
        \hline
        F & T & F \\
        \hline
        F & F & T \\
        \hline
        \end{tabular}

        This is logically equivalent to $(p\Rightarrow q )\land (q\Rightarrow p)$. In
        words this is often written ``$p$ if and only if $q$".
        \end{definition}
}{% presentation
}{% nopresentation
}{% comments
}

\content{
        \begin{example} Suppose that this statement is true: ``The garbage truck
        comes down my street if and only if it is Tuesday morning." Which of the
        following statements could be true? 
        \begin{enumerate}
        \item It is noon on Tuesday, and the garbage truck did not come down my
        street this morning.
        \item It is Monday, and the garbage truck came down my street this
        morning.
        \item It is Wednesday, and the garbage truck did not come down my street
        this morning.
        \end{enumerate}
        \end{example}
}{% presentation
}{% nopresentation
}{% comments
}
